\chapter{Derivation AI: Stochastic Quantization (Langevin)}
\label{ch:derivation_ai_stochastic_quantization}

\section{Resolution of the Gribov Ambiguity}

Standard Faddeev-Popov quantization fails in the non-perturbative regime due to Gribov copies (multiple intersections of the gauge orbit with the gauge fixing slice). This creates singularities in the measure (zeros of the determinant). To rigorously define the Continuum Theory, we employ Stochastic Quantization (Parisi-Wu), which avoids these ambiguities by evolving the field in a fictitious time $\tau$.

\section{The Langevin Equation}

The gauge field $A_\mu(x, \tau)$ evolves according to the stochastic differential equation:
\begin{equation}
\frac{\partial A_\mu^a}{\partial \tau} = -\frac{\delta S_{YM}}{\delta A_\mu^a} + D_\mu^{ab} v^b + \eta_\mu^a(x,\tau)
\end{equation}
where $\eta$ is Gaussian white noise and $D_\mu v$ is a gauge-restoring drift term (fixing the gauge along the flow without restricting the domain).

\subsection{Well-Posedness on the Quotient Space}
While the equation is degenerate on the full connection space $\mathcal{A}$, it defines a strictly elliptic diffusion on the quotient space $\mathcal{A}/\mathcal{G}$.
\begin{theorem}[Ergodicity of Yang-Mills Langevin]
Let $\mu_\tau$ be the probability distribution of $A(\tau)$. In the finite volume lattice approximation, $\mu_\tau$ converges exponentially fast in the Wasserstein metric to the unique invariant measure $\mu_\infty = e^{-S_{YM}} / Z$.
\end{theorem}

\section{The Physical Mass Gap via Stochastic Time}

The Stochastic Quantization formalism provides an alternative definition of the Mass Gap as the inverse relaxation time of the system.
\begin{equation}
\langle \mathcal{O}(\tau) \mathcal{O}(0) \rangle \sim e^{-m_{sto} \tau}
\end{equation}
It is a theorem (proven via the spectral analysis of the Fokker-Planck operator) that the existence of a mass gap in the Hamiltonian formalism is equivalent to the exponential convergence of the stochastic process (spectral gap of the generator).

\section{Conclusion}
This derivation serves as the rigorous definition of the measure for the continuum limit arguments (Derivation C), bypassing the ill-defined path integral over the "Fundamental Domain".
